\documentclass[11pt]{article}
\usepackage{fontspec}
\setmainfont{Times New Roman}
\setsansfont{Arial}
\setmonofont{Consolas}
\usepackage{amsmath,amssymb,mathtools}
\usepackage{geometry}
\usepackage{microtype}
\geometry{a4paper,margin=2.35cm}
\setlength{\parindent}{1.25em}
\setlength{\parskip}{0pt}
\makeatletter
\renewcommand\footnotesize{\@setfontsize\footnotesize{7.7}{8.7}\selectfont}
\makeatother
\emergencystretch=65em
\allowdisplaybreaks
\sloppy
\hbadness=10000
\vbadness=10000
\hfuzz=2pt
\vfuzz=10pt
\raggedbottom

\providecommand{\mR}{\mathfrak{R}}
\providecommand{\ma}{\mathfrak{a}}
\providecommand{\mb}{\mathfrak{b}}
\providecommand{\mc}{\mathfrak{c}}
\providecommand{\mm}{\mathfrak{m}}
\providecommand{\mn}{\mathfrak{n}}
\providecommand{\mo}{\mathfrak{o}}
\providecommand{\mq}{\mathfrak{q}}
\providecommand{\mt}{\mathfrak{t}}
\providecommand{\mpideal}{\mathfrak{p}}
\providecommand{\mbar}[1]{\overline{#1}}

\begin{document}

% Unit: Paper30 Section08 v001. Direct Interslavic Latin translation from the German final-audited source slice, source lines 733--858 / segments P30-S0130--P30-S0140. Footnote counter continues after Paper30 Section07.
\setcounter{footnote}{29}

\subsection*{\S~8. Teorija idealov pri integralnoj zamknjenosti v polju kvocientov}

Do sego razvita teorija idealov teper ima byti zaoštrjena do obyčnoj formy dodanjem dvuh poslědnjih aksiomov.

1. \emph{Pomočny teorem.} Naj $\mR$ bude komutativno kolco bez nulovyh děliteljev s jediničnym elementom, i naj v $\mR$ bude predpostavljeno uslovje cěpov děliteljev (aksiomy I, III, IV). Togda iz $\ma=\ma\mb$ i $\ma\ne(0)$ vsegda slěduje $\mb=\mo$. Znači,
\[
        \mc\ne\mc\mt
\]
za vse idealy različne od nulovogo i jediničnogo idealov.
\footnote{Protivno tomu, pri tyh samyh predpostavkah iz $\ma\mb=\ma\mc$ ne možno zaključiti $\mb=\mc$. Kako primjer možno vzęti kolco polinomov $K[x,y]$: $\ma=(xy)$, $\mb=(x^3,xy,y^2)$, $\mc=(x^2,y^2)$; dejansko $\ma\mb=\ma\mc=\ma^2$, ale $\mb\ne\mc$.}
Dokaz vyhodi kako pri pomočnom teoremu §1, 3, bo $\ma\mb$ možno razsmatrjati kako konečny modul vzhodno k $\mb$. Ako $\alpha_1,\ldots,\alpha_n$ označa bazu ideala $\ma$, ktora obstaja po I, togda iz $\ma=\ma\mb$ slěduje sistem jednačenj
\[
        \alpha_i=b_{i1}\alpha_1+\cdots+b_{in}\alpha_n,
        \qquad b_{ij}\equiv0\pmod{\mb},\quad i=1,\ldots,n.
\]
Poněže $\mR$ jest predpostavljeno bez nulovyh děliteljev, iz togo slěduje $|b_{ij}-e_{ij}|=0$, kde $e_{ij}=0$ za $i\ne j$, a $e_{ii}=e$. Iz jednačenja
\[
        e-b_{11}-\cdots\equiv0\pmod{\mb}
\]
slěduje ale $e\equiv0\pmod{\mb}$, i zato $\mb=\mo$.

2. Teper ostava samo dokazati, že po dodanju predpostavky integralnoj zamknjenosti v polju kvocientov (aksioma V) k aksiomam I--IV vsaky primarny ideal različny od nulovogo ideala stava stepenju svojego asociovanogo prostogo ideala. Jednoznačnost slědujučego predstavjenja
\[
        \mm=\mpideal_1^{\varrho_1}\cdots\mpideal_r^{\varrho_r}
\]
za vse idealy različne od nulovogo ideala uže jest dokazana teoremoju V i pomočnym teoremom pod 1; jednočasno jest pokazano, že ty proste idealy ne imajut vlastnogo dělitelja različnogo od jediničnogo ideala. Dokaz bude najprvo proveden za eksponent dva s upotrěboju Dedekindovogo poslědka II (§1, 4), a v obče črez polnu indukciju.

\emph{Pomočny teorem.} \emph{Pri predpostavce aksiomov I--V ne obstajut drugi primarne idealy eksponenta dva osim $\mq=\mpideal^2$.}

Treba dokazati, že iz
\[
        \mpideal^2\equiv0\pmod{\mq},
        \qquad \mq\equiv0\pmod{\mpideal},
        \qquad \mq\not\equiv0\pmod{\mpideal^2}
\]
nužno slěduje $\mq=\mpideal^2$. Naj zato bude
\[
        \mc\equiv0\pmod{\mq},
        \qquad \mc\not\equiv0\pmod{\mpideal},
        \qquad \mc\equiv0\pmod{\mpideal^2}.
\]
Togda iz pomočnogo teorema §7, 3 slěduje, že $\mc:\mpideal$ stava vlastnym děliteljem $\mc$. Zato obstaja element $b$ takovy, že
\[
        b\mpideal\equiv0\pmod{\mc}\equiv0\pmod{\mq},
        \qquad b\not\equiv0\pmod{\mc}.
\]
Tym $y=b/c$ ne jest cěly, to jest $y$ naleži k polju kvocientov, ale ne k $\mR$. Treba dokazati $b\not\equiv0\pmod{\mpideal}$; iz togo slěduje, poněže $\mq$ jest primarny i naleži k $\mpideal$, že $\mpideal\equiv0\pmod{\mq}$.

Po Dedekindovom poslědku II obstajut elementy $m,n$ iz $\mR$ take, že
\[
        y=b/c=m/n
\]
i takodže $m^2/n$ ne jest cěly, to jest
\[
        bn=mc,
        \qquad m\not\equiv0\pmod{\mn},
        \qquad m^2\not\equiv0\pmod{\mn}.
\]
Množenje $b\mpideal$ s $m$ daje
\[
        mb\mpideal\equiv0\pmod{\mn c},
        \qquad\text{zato}\qquad m\mpideal\equiv0\pmod{\mn},
\]
poslědnje, bo idě o kolco bez nulovyh děliteljev. Iz togo slěduje
\[
        m\not\equiv0\pmod{\mpideal},
\]
zbog $m^2\not\equiv0\pmod{\mn}$ i $\mn\not\equiv0\pmod{\mpideal}$; poslědnje znovu po pomočnom teoremu §7, 3, bo $\mn:\mpideal$ jest vlastny dělitelj $\mn$. Četyri relacije
\[
        m\not\equiv0\pmod{\mpideal},
        \quad n\not\equiv0\pmod{\mpideal},
        \quad c\equiv0\pmod{\mpideal},
        \quad c\not\equiv0\pmod{\mpideal^2},
        \quad bn=mc
\]
davajut $b\not\equiv0\pmod{\mpideal}$. Dejateljno, poněže $\mpideal^2$ jest primarny, najprvo dobivamo $mc\not\equiv0\pmod{\mpideal^2}$, i zato iz $bn=mc$ slěduje $b\not\equiv0\pmod{\mpideal}$. Tym pomočny teorem jest dokazany.

3. \emph{Pomočny teorem.} \emph{Pri predpostavce aksiomov I--V ne obstajut drugi primarne idealy eksponenta $\varrho$ osim $\mq=\mpideal^\varrho$.}
\footnote{Dokaz pomočnogo teorema 3 pri predpostavce validnosti pomočnogo teorema 2 nahodi se u Masazo Sono, \emph{On Congruences II}, §§9 i 10.}
Pomočny teorem slěduje iz pomočnogo teorema pod 2 bez daljnjego upotrěbljenja aksiomov. Za dokaz treba najprvo pokazati:
\[
        \text{ako }\mc\equiv0\pmod{\mpideal},\ \mc\not\equiv0\pmod{\mpideal^2},
        \quad\text{togda}\quad
        \mpideal^\varrho=(\mc^\varrho,\mpideal^{\varrho+1})
\]
za vsako $\varrho$. Po pomočnom teoremu pod 2 imamo $\mpideal=(\mc,\mpideal^2)$. Ako zato predpostavimo
\[
        \mpideal^{\varrho-1}=(\mc^{\varrho-1},\mpideal^\varrho),
\]
togda črez množenje s $\mpideal$ po distributivnom zakonu slěduje
\[
        \mpideal^\varrho=(\mc^\varrho,\mpideal^{\varrho+1}).
\]
Dokazujemny pomočny teorem možno privesti k slědujučej drugoj formě:
\[
        \mq\equiv0\pmod{\mpideal^\varrho},
        \quad \mq\not\equiv0\pmod{\mpideal^{\varrho+1}},
        \quad \mpideal^{\varrho+1}\equiv0\pmod{\mq},
        \quad \varrho\ge1
\]
implikuje
\[
        \mpideal^\varrho\equiv0\pmod{\mq}.
\]
Bo zaradi $\mq\equiv0\pmod{\mpideal^\varrho}$ i $\mpideal^\varrho\ne\mpideal^{\varrho+1}$ po pomočnom teoremu pod 1 vtoro uslovje jest ispolnjeno. Konečno kratno upotrěbljenje toj drugoj formy daje $\mpideal^\varrho\equiv0\pmod{\mq}$ i s tym $\mq=\mpideal^\varrho$.

Iz predpostavky drugoj formy, s ohledom na predstavjenje $\mpideal^\varrho$, za vsaky element $q$ iz $\mq$ slěduje:
\[
        q=b c^\varrho+p^{\varrho+1},
        \qquad b\equiv0\pmod{\mpideal},
\]
ili, ravnoznačno,
\[
        bc^\varrho\equiv0\pmod{\mq,\mpideal^{\varrho+1}}.
\]
Poněže $(\mq,\mpideal^{\varrho+1})$ jest primarny, iz togo slěduje
\[
        c^\varrho\equiv0\pmod{\mq,\mpideal^{\varrho+1}}
        \quad\text{ili}\quad
        c^{\varrho+1}\equiv0\pmod{\mq,\mpideal^{\varrho+2}},
\]
i s tym $\mpideal^\varrho\equiv0\pmod{\mq}$.

4. Sumarno dobivamo:

\textbf{Teorema VI.} \emph{Ako v komutativnom kolcu važe aksiomy I--V, togda vsaky ideal različny od nulovogo i jediničnogo idealov možno jednoznačno predstaviti kako produkt stepenjev konečno mnogih prostyh idealov različnyh od nulovogo i jediničnogo idealov; ty proste idealy ne imajut vlastnogo dělitelja različnogo od jediničnogo ideala.}

\end{document}
